\documentclass[10pt,a4paper]{amsart}
\usepackage[margin=19mm]{geometry}
\usepackage{amsmath,amssymb,mathtools}
\usepackage[scr=rsfs,cal=euler]{mathalfa}
\usepackage{xfrac}
\usepackage[unicode,hidelinks,bookmarks=false]{hyperref}
\newcommand{\ie}{i.e.\ }
\newcommand{\cf}{cf.\ }
\newcommand{\resp}{respectively\ }
\newcommand{\Subsection}{\subsection}
\providecommand{\nobreakdash}{\nobreak\hbox{-}}
\setlength{\emergencystretch}{2em}
\multlinegap0pt
\usepackage{amssymb}
\usepackage{mathtools}
\usepackage[shortlabels]{enumitem}
\newtheorem{lemma}{Lemma}
\newtheorem{proposition}{Proposition}
\newcommand{\NN}{\mathbb N}
\newcommand{\RR}{\overline{\mathbb R}_{+}}
\newcommand{\up}{\mathord\uparrow}
\title{A continuous $3$-distributive frame that is not $\omega$-distributive}
\author{Wei Luan, Qingguo Li}
\date{Source: arXiv:2609.11671v1 (CC BY 4.0)}
\begin{document}
\maketitle

\noindent\emph{Source and licence.} A continuous $3$-distributive frame that is not $\omega$-distributive. Version arXiv:2609.11671v1 (CC BY 4.0), distributed by arXiv under the Creative Commons Attribution 4.0 International licence, http://creativecommons.org/licenses/by/4.0/, as recorded in the arXiv licence metadata for this exact version and shown on the abstract page https://arxiv.org/abs/2609.11671v1. Full licence notice: \texttt{licenses/arxiv-2609.11671-CC-BY-4.0.txt}.\par\smallskip

\noindent\textit{Editorial scope.} Counted unit: the article's proposition prop:continuous-meet (source0.tex lines 496--499). The article's complete original proof of it is delivered with it (source0.tex lines 501--541, one \texttt{proof} environment of 41 lines). No mathematics is written, repaired or re-derived: the statement and the proof are the article's own lines, in the article's own notation, and no retained line is substituted or re-flowed.  Retained uncounted as the source-local context the proof needs: lines 304--312; lines 390--393.  Nothing was omitted from any retained span, so the package carries no omission record.  No retained line contains a citation, so the wrapper carries no bibliography and no bibliography entry.  No retained line contains a figure environment, an image-inclusion command or drawing code, so no image is delivered and the package declares no asset dependency.  Editorial formatting only, in addition: an AMS \texttt{amsart} preamble that restates the article's own declarations and macros used by the retained lines, namely the environments \texttt{lemma}, \texttt{proposition} on separate counters, together with the standard packages those lines need.  The retained environments print in this document as Lemma 1, Proposition 1 in order of appearance, whereas the article prints the same items with the numbering of its own full document.  The complete native source of the article as distributed in the arXiv e-print for this exact version is bundled unchanged as \texttt{sources/210043/source0.tex}.
\begin{lemma}\label{lem:slow-intersection}
For real numbers $\tau>\varepsilon>0$, there is $i_0\in\NN_0$ such
that for all $i\geq i_0$ and $A,B\in S_i$,
\[
\sigma_i(A)>\tau,\quad\sigma_i(B)>\tau
\quad\Longrightarrow\quad
\sigma_i(A\cap B)>\tau-\varepsilon.
\]
\end{lemma}

\begin{equation}\label{eq:meet}
x\wedge y=(u,(A_i\cap B_i)),\qquad
u=\min\left\{s,t,\inf_{i\in\NN_0}\sigma_i(A_i\cap B_i)\right\}.
\end{equation}

\begin{proposition}\label{prop:continuous-meet}
The map $\wedge:X\times X\longrightarrow X$ in \eqref{eq:meet}
is continuous.
\end{proposition}

\begin{proof}
Let $x=(s,(A_i))$, $y=(t,(B_i))$, and suppose
$x\wedge y\in W=B(J;m,E)$. Write the first coordinate of
$x\wedge y$ as $u$. For $i\leq m$, we have $E_i\subseteq A_i\cap B_i$.
If $J=\RR$, then $x,y\in W$ and every meet of two points of $W$
belongs to $W$. Hence $W\times W\subseteq\wedge^{-1}(W)$.

Suppose $J=(r,\infty]$. Since $u>r$, choose finite real numbers
$\tau,\varepsilon$ such that
\begin{equation}\label{eq:margin}
r<\tau-\varepsilon<\tau<u,\qquad\varepsilon>0.
\end{equation}
Then $\tau>\varepsilon$, $s,t>\tau$, and
$\sigma_i(A_i\cap B_i)>\tau$ for every $i$.
Choose $i_0$ from Lemma~\ref{lem:slow-intersection}, and put
$q=\max\{m,i_0\}$. The sets
\begin{align*}
P&=X\cap\left((\tau,\infty]\times\prod_{i=0}^{q}\up A_i
                                      \times\prod_{i>q}S_i\right),\\
Q&=X\cap\left((\tau,\infty]\times\prod_{i=0}^{q}\up B_i
                                      \times\prod_{i>q}S_i\right)
\end{align*}
are open neighborhoods of $x$ and $y$.
Take $x'=(s',(A_i'))\in P$ and $y'=(t',(B_i'))\in Q$.
For $i\leq q$, monotonicity gives
\[
\sigma_i(A_i'\cap B_i')\geq\sigma_i(A_i\cap B_i)>\tau.
\]
For $i>q$, the defining constraints of $X$ imply
$\sigma_i(A_i')\geq s'>\tau$ and
$\sigma_i(B_i')\geq t'>\tau$. Lemma~\ref{lem:slow-intersection}
therefore gives $\sigma_i(A_i'\cap B_i')>\tau-\varepsilon$.
Consequently
\[
\inf_i\sigma_i(A_i'\cap B_i')\geq\tau-\varepsilon,
\]
so the first coordinate of $x'\wedge y'$ is at least
$\tau-\varepsilon>r$. Also $E_i\subseteq A_i'\cap B_i'$ for
$i\leq m$. Hence $x'\wedge y'\in W$, proving
$P\times Q\subseteq\wedge^{-1}(W)$.
\end{proof}

\end{document}
